\documentclass[12pt]{article}
% Спільна преамбула для звітів (XeLaTeX)
\usepackage{fontspec}
\usepackage{polyglossia}
\setmainlanguage{ukrainian}
\setotherlanguage{english}
\setmainfont{DejaVuSerif}[Extension=.ttf,UprightFont=*,BoldFont=*-Bold,ItalicFont=*-Italic,BoldItalicFont=*-BoldItalic]
\setsansfont{DejaVuSans}[Extension=.ttf,UprightFont=*,BoldFont=*-Bold,ItalicFont=*-Oblique,BoldItalicFont=*-BoldOblique]
\setmonofont{DejaVuSansMono}[Extension=.ttf,UprightFont=*,BoldFont=*-Bold,ItalicFont=*-Oblique,BoldItalicFont=*-BoldOblique,Scale=0.9]
\usepackage[a4paper,margin=2.2cm]{geometry}
\usepackage{booktabs}
\usepackage{tabularx}
\usepackage{array}
\usepackage{enumitem}
\usepackage{xcolor}
\usepackage{tcolorbox}
\usepackage{graphicx}
\usepackage{setspace}
\usepackage[hidelinks]{hyperref}
\usepackage{titlesec}
\usepackage{fancyhdr}
\usepackage{multirow}
\usepackage{colortbl}
\setlength{\parskip}{4pt}
\setlength{\parindent}{0pt}
\onehalfspacing

\definecolor{accent}{HTML}{1F4E79}
\definecolor{sgreen}{HTML}{E2F0D9}
\definecolor{wred}{HTML}{FBE5D6}
\definecolor{oblue}{HTML}{DEEBF7}
\definecolor{tyellow}{HTML}{FFF2CC}
\definecolor{lightgray}{HTML}{F2F2F2}

\titleformat{\section}{\large\bfseries\color{accent}}{\thesection}{0.6em}{}
\titleformat{\subsection}{\normalsize\bfseries}{\thesubsection}{0.6em}{}

\pagestyle{fancy}
\fancyhf{}
\lhead{\small ІТ, Модуль 2 — командний проєкт}
\rhead{\small Петров Б., Мекшаков А.}
\cfoot{\thepage}
\renewcommand{\headrulewidth}{0.3pt}

\newcolumntype{L}[1]{>{\raggedright\arraybackslash}p{#1}}
\newcolumntype{Y}{>{\raggedright\arraybackslash}X}

% Титульна сторінка
\newcommand{\maketitlepage}[3]{%
\begin{titlepage}
\centering
{\large Київський національний університет імені Тараса Шевченка\\[2pt]
Факультет комп'ютерних наук та кібернетики\par}
\vspace{3cm}
{\large Дисципліна: \textbf{Інформаційні технології}\par}
\vspace{0.5cm}
{\large Модуль 2 — робота в команді\par}
\vspace{2cm}
{\LARGE\bfseries #1\par}
\vspace{0.6cm}
{\large #2\par}
\vfill
\begin{flushright}
\begin{tabular}{l}
\textbf{Виконали:} студенти 4 курсу, група МІ-41\\[4pt]
Петров Богдан\\
Мекшаков Артемій\\[10pt]
\textbf{Спеціальність:} Комп'ютерні науки\\
\end{tabular}
\end{flushright}
\vspace{1.5cm}
{Київ — #3\par}
\end{titlepage}
}

\usepackage{pdflscape}
\usepackage{longtable}
\usepackage{pgfgantt}

\begin{document}
\maketitlepage{Звіт за етапом 3}{Планування проєкту: дорожня карта, RACI-матриця, пріоритезація беклогу, обсяг семестру}{2026}

\section{Мета етапу}
Спланувати життєвий цикл проєкту StudyFlow: описати проєкт (Project Overview), побудувати діаграму Ґанта на один рік, розподілити відповідальність між учасниками (RACI-матриця), сформувати і пріоритезувати беклог продукту (матриця Ейзенхауера, MoSCoW) та визначити підмножину функціоналу для реалізації в межах модуля~2.

\section{Дорожня карта проєкту}
\subsection{Project Overview}
\begin{tabularx}{\textwidth}{L{3.4cm} Y}
\toprule
\textbf{Назва} & StudyFlow --- планувальник навчальних дедлайнів для студентських груп\\
\midrule
\textbf{Мета} & Дати студентській групі єдине місце для дедлайнів, задач і календаря, щоб зменшити кількість пропущених термінів і ручне дублювання інформації старостою\\
\midrule
\textbf{Цілі (SMART)} &
1) До 18.11.2026 --- робочий прототип з автентифікацією, групами, дедлайнами та календарем тижня, покритий $\geq$20 модульними тестами;\newline
2) До 01.12.2026 --- demo-версія з канбаном, нагадуваннями й прогресом групи, фінальний звіт;\newline
3) До 02.2027 --- пілот у 2--3 групах, $\geq$70\,\% дедлайнів групи заведено в системі;\newline
4) До 10.2027 --- версія 2.5, готова до впровадження на рівні факультету\\
\midrule
\textbf{Обсяг (in scope)} & Вебзастосунок (Next.js + NestJS + PostgreSQL): користувачі, групи та ролі, предмети, особисті й групові дедлайни, календар, канбан задач, e-mail-нагадування, прогрес групи; API-документація; тести; деплой\\
\midrule
\textbf{Поза обсягом семестру} & Інтеграції з Google Classroom / Calendar, мобільний застосунок / PWA, Telegram-бот, адмін-панель факультету, монетизація\\
\midrule
\textbf{Очікувані результати} & Репозиторій з кодом та CI; задокументоване REST API (Swagger); UML-моделі та схема БД; журнали модульних та інтеграційних тестів; публічна demo-версія; звіти за етапами й фінальна презентація\\
\midrule
\textbf{Зацікавлені сторони} & Команда (2 розробники); викладач дисципліни (замовник, приймає результати етапів); студенти групи МІ-41 (ранні користувачі)\\
\midrule
\textbf{Обмеження} & Команда з 2 осіб; жорсткі дедлайни етапів (03.11, 18.11, 01.12); бюджет 0 грн --- лише безкоштовні тарифи хостингу та e-mail; паралельне навантаження з іншими дисциплінами\\
\midrule
\textbf{Припущення} & Стек відомий обом учасникам; одногрупники погодяться на пілот; безкоштовних лімітів сервісів достатньо для пілоту\\
\midrule
\textbf{Ризики} & Не встигнути з обсягом MVP (мітигація: пріоритезація MoSCoW, Could-функції переносяться); хвороба/недоступність учасника (мітигація: обидва знають весь стек, рев'ю коду); проблеми з доставкою e-mail (мітигація: ізольований модуль нагадувань, запасний провайдер)\\
\bottomrule
\end{tabularx}

\begin{landscape}
\subsection{Діаграма Ґанта на рік}
Діаграма охоплює період з листопада 2026 до жовтня 2027 включно --- повний життєвий цикл проєкту відповідно до карти продукту, незалежно від меж семестру. Етапи модуля~2 (ініціація, планування, прототип, контроль) виділено як перші чотири блоки.

\begin{center}
\resizebox{\linewidth}{!}{%
\begin{ganttchart}[
  x unit=1.65cm, y unit title=0.7cm, y unit chart=0.55cm,
  hgrid, vgrid={*1{dotted}},
  title/.style={fill=accent, draw=accent}, title label font=\bfseries\color{white}\footnotesize,
  bar/.style={fill=oblue, draw=accent!60!black, rounded corners=1pt},
  bar label font=\footnotesize, bar height=0.6,
  group/.style={fill=accent, draw=accent}, group label font=\bfseries\footnotesize,
  group height=0.35, group peaks tip position=0,
  milestone/.style={fill=wred!80!red, draw=red!60!black}, milestone label font=\footnotesize,
  link/.style={-latex, accent!70!black, thick}
]{1}{12}
\gantttitle{Листопад 2026 --- Жовтень 2027}{12}\\
\gantttitlelist{11,12,1,2,3,4,5,6,7,8,9,10}{1}\\
\ganttgroup{Модуль 2: ініціація, планування, прототип, звіт}{1}{1}\\
\ganttbar{Етапи 2--3: Lean Canvas, roadmap, RACI, беклог}{1}{1}\\
\ganttbar{Етап 4: MVP 0.1 (спринти 1--3)}{1}{1}\\
\ganttbar{Етап 5: версія 0.2, фінальний звіт}{1}{2}\\
\ganttmilestone{Захист модуля 2 (01.12)}{2}\\
\ganttgroup{Версія 1.0: пілот і якість}{2}{4}\\
\ganttbar{Пілот у 2--3 групах, збір зворотного зв'язку}{2}{4}\\
\ganttbar{Виправлення, відновлення пароля, налаштування}{3}{4}\\
\ganttbar{Особиста статистика, експорт iCal}{4}{4}\\
\ganttmilestone{Реліз 1.0}{4}\\
\ganttgroup{Версія 1.5: інтеграції}{5}{7}\\
\ganttbar{Імпорт із Google Classroom}{5}{6}\\
\ganttbar{Синхронізація з Google Calendar}{6}{7}\\
\ganttbar{Дайджест, коментарі до задач}{7}{7}\\
\ganttmilestone{Реліз 1.5}{7}\\
\ganttgroup{Версія 2.0: мобільність та канали}{8}{10}\\
\ganttbar{PWA, push-сповіщення}{8}{9}\\
\ganttbar{Telegram-бот}{9}{10}\\
\ganttbar{Графік навантаження групи}{10}{10}\\
\ganttmilestone{Реліз 2.0}{10}\\
\ganttgroup{Версія 2.5: масштабування}{11}{12}\\
\ganttbar{Модель факультету, адмін-панель}{11}{12}\\
\ganttbar{Аналітика для викладачів}{12}{12}\\
\ganttbar{Ретроспектива року, план на наступний рік}{12}{12}\\
\ganttmilestone{Реліз 2.5}{12}
\end{ganttchart}}
\end{center}
\end{landscape}

\subsection{RACI-матриця}
Позначення: \textbf{R} (Responsible) --- виконує; \textbf{A} (Accountable) --- відповідає за результат, приймає; \textbf{C} (Consulted) --- консультує; \textbf{I} (Informed) --- інформується. Для кожного завдання рівно один~A.

\begin{table}[h]
\centering
\renewcommand{\arraystretch}{1.2}
\begin{tabularx}{\textwidth}{Y c c c}
\toprule
\textbf{Завдання / зона відповідальності} & \textbf{Петров Б.} & \textbf{Мекшаков А.} & \textbf{Викладач}\\
\midrule
Управління проєктом, графік, контроль дедлайнів & \textbf{A}/R & C & I\\
Вимоги, пріоритезація беклогу, ведення дошки & C & \textbf{A}/R & C\\
Архітектура системи, вибір інструментарію & \textbf{A}/R & C & I\\
Проєктування схеми БД, міграції & \textbf{A}/R & C & ---\\
Backend: REST API, бізнес-логіка, автентифікація & \textbf{A}/R & I & ---\\
Модуль нагадувань (cron, e-mail) & \textbf{A}/R & I & ---\\
UI/UX: макети, дизайн-система & I & \textbf{A}/R & ---\\
Frontend: сторінки, компоненти, інтеграція з API & C & \textbf{A}/R & ---\\
Календар та канбан-дошка (клієнт) & C & \textbf{A}/R & ---\\
Модульні тести backend, журнал тестування & \textbf{A}/R & C & I\\
Інтеграційні та E2E-тести & C & \textbf{A}/R & I\\
CI/CD, Docker, деплой demo-версії & \textbf{A}/R & I & ---\\
API-документація (Swagger) & \textbf{A}/R & C & I\\
UML-моделювання & R & \textbf{A}/R & I\\
Код-рев'ю pull request'ів & R & R & ---\\
Звіти за етапами, презентації, демо & R & \textbf{A}/R & C\\
Пілот у групі, збір зворотного зв'язку & C & \textbf{A}/R & I\\
Приймання результатів етапів & I & I & \textbf{A}\\
\bottomrule
\end{tabularx}
\caption{RACI-матриця ролей}
\end{table}

\section{Пріоритезація функціоналу}
Беклог сформовано з 45 елементів (33 user story, 9 технічних і 3 організаційні задачі), об'єднаних у 9 епіків. Кожен елемент оцінено у story points (шкала 1--2--3--5--8), пріоритезовано за MoSCoW та класифіковано за матрицею Ейзенхауера.

\subsection{Матриця Ейзенхауера}
Критерій \emph{важливості} --- внесок у ціннісну пропозицію та вимоги курсу; критерій \emph{терміновості} --- необхідність для найближчого демо або дедлайну етапу.

\renewcommand{\arraystretch}{1.1}
\footnotesize
\begin{tabularx}{\textwidth}{|Y|Y|}
\hline
\cellcolor{wred}\textbf{Q1. Важливо + терміново --- робити зараз (спринти 1--2)} & \cellcolor{oblue}\textbf{Q2. Важливо, не терміново --- планувати (спринти 3--4 і далі)}\\
\hline
\begin{itemize}[leftmargin=1em,nosep]\item T-01 Монорепозиторій Next.js + NestJS, лінтери, Prettier
\item T-02 Docker Compose: PostgreSQL + backend + frontend
\item T-04 Базовий UI-кіт та макет сторінок
\item US-01 Реєстрація за e-mail і паролем
\item US-02 Вхід / вихід, JWT-сесія
\item US-03 Сторінки реєстрації та входу
\item US-06 Створення навчальної групи
\item US-07 Приєднання за інвайт-кодом
\item US-08 Ролі: староста / учасник
\item US-10 CRUD предметів групи
\item US-11 CRUD особистих дедлайнів
\item US-12 Групові дедлайни (створює староста, бачать усі)
\item US-13 Список дедлайнів з фільтрами й сортуванням
\item US-14 Відмітка виконання дедлайну
\item US-15 Форми створення/редагування дедлайну
\item US-16 Календар на тиждень\end{itemize} &
\begin{itemize}[leftmargin=1em,nosep]\item T-03 CI (GitHub Actions): лінт, збірка, тести
\item T-05 Swagger / OpenAPI-документація API
\item T-06 Модульні тести backend ($\geq$20)
\item T-07 Інтеграційні та E2E-тести
\item T-08 UML-діаграми та схема БД
\item T-09 Деплой на хостинг (Vercel + Railway/Render)
\item US-04 Профіль користувача (ім'я, група, зміна пароля)
\item US-09 Сторінка групи: учасники, керування
\item US-17 Календар на місяць
\item US-18 Експорт у iCal (.ics)
\item US-19 Канбан-дошка командних задач (To Do / Doing / Done)
\item US-20 API задач: CRUD, статуси, виконавець
\item US-21 Drag-and-drop карток між колонками
\item US-23 E-mail-нагадування за N днів до дедлайну
\item US-24 Налаштування нагадувань користувачем
\item US-27 Прогрес групи по дедлайну (\% виконали)
\item US-30 Імпорт дедлайнів з Google Classroom
\item US-31 Синхронізація з Google Calendar
\item D-03 Фінальний звіт за модуль 2 і фінальна презентація\end{itemize}\\
\hline
\cellcolor{tyellow}\textbf{Q3. Терміново, не важливо --- мінімізувати зусилля} & \cellcolor{lightgray}\textbf{Q4. Не важливо, не терміново --- після семестру}\\
\hline
\begin{itemize}[leftmargin=1em,nosep]\item D-01 Звіти за етапами 2–4 у класрум
\item D-02 Щотижневі демо-презентації та рефлексія спринта\end{itemize} &
\begin{itemize}[leftmargin=1em,nosep]\item US-05 Відновлення пароля через e-mail
\item US-22 Коментарі до задач
\item US-25 Щотижневий дайджест на e-mail
\item US-26 Telegram-бот для нагадувань
\item US-28 Особиста статистика виконання
\item US-29 Навантаження групи по тижнях (графік)
\item US-32 PWA / мобільна версія
\item US-33 Масштабування на факультет: кілька груп, адмін-панель\end{itemize}\\
\hline
\end{tabularx}
\normalsize
\renewcommand{\arraystretch}{1.3}

\subsection{Беклог продукту}
Пріоритети MoSCoW: \textbf{Must} --- без цього прототип не має сенсу; \textbf{Should} --- важливо, але можна відкласти в межах семестру; \textbf{Could} --- бажано за наявності часу; \textbf{Won't} --- свідомо поза семестром. Колонка Q --- квадрант матриці Ейзенхауера.

\begin{landscape}
\footnotesize
\renewcommand{\arraystretch}{1.15}
\begin{longtable}{l L{7.2cm} L{3.4cm} l c l l c}
\toprule
\textbf{ID} & \textbf{Назва} & \textbf{Епік} & \textbf{MoSCoW} & \textbf{Q} & \textbf{Спринт} & \textbf{Виконавець} & \textbf{SP}\\
\midrule
\endfirsthead
\toprule
\textbf{ID} & \textbf{Назва} & \textbf{Епік} & \textbf{MoSCoW} & \textbf{Q} & \textbf{Спринт} & \textbf{Виконавець} & \textbf{SP}\\
\midrule
\endhead
\bottomrule
\endfoot
T-01 & Монорепозиторій Next.js + NestJS, лінтери, Prettier & Інфраструктура та якість & Must & Q1 & Sprint 1 & Богдан & 2\\
T-02 & Docker Compose: PostgreSQL + backend + frontend & Інфраструктура та якість & Must & Q1 & Sprint 1 & Богдан & 2\\
T-03 & CI (GitHub Actions): лінт, збірка, тести & Інфраструктура та якість & Should & Q2 & Sprint 1 & Богдан & 2\\
T-04 & Базовий UI-кіт та макет сторінок & Інфраструктура та якість & Must & Q1 & Sprint 1 & Артемій & 3\\
T-05 & Swagger / OpenAPI-документація API & Інфраструктура та якість & Must & Q2 & Sprint 3 & Богдан & 2\\
T-06 & Модульні тести backend ($\geq$20) & Інфраструктура та якість & Must & Q2 & Sprint 3 & Богдан & 3\\
T-07 & Інтеграційні та E2E-тести & Інфраструктура та якість & Must & Q2 & Sprint 3 & Артемій & 3\\
T-08 & UML-діаграми та схема БД & Інфраструктура та якість & Must & Q2 & Sprint 2 & Обидва & 2\\
T-09 & Деплой на хостинг (Vercel + Railway/Render) & Інфраструктура та якість & Could & Q2 & Sprint 4 & Богдан & 2\\
US-01 & Реєстрація за e-mail і паролем & Автентифікація & Must & Q1 & Sprint 1 & Богдан & 2\\
US-02 & Вхід / вихід, JWT-сесія & Автентифікація & Must & Q1 & Sprint 1 & Богдан & 2\\
US-03 & Сторінки реєстрації та входу & Автентифікація & Must & Q1 & Sprint 1 & Артемій & 2\\
US-04 & Профіль користувача (ім'я, група, зміна пароля) & Автентифікація & Should & Q2 & Sprint 2 & Артемій & 2\\
US-05 & Відновлення пароля через e-mail & Автентифікація & Could & Q4 & Після семестру & Богдан & 2\\
US-06 & Створення навчальної групи & Групи & Must & Q1 & Sprint 1 & Богдан & 2\\
US-07 & Приєднання за інвайт-кодом & Групи & Must & Q1 & Sprint 1 & Обидва & 2\\
US-08 & Ролі: староста / учасник & Групи & Must & Q1 & Sprint 1 & Богдан & 2\\
US-09 & Сторінка групи: учасники, керування & Групи & Should & Q2 & Sprint 2 & Артемій & 2\\
US-10 & CRUD предметів групи & Предмети та дедлайни & Must & Q1 & Sprint 2 & Богдан & 2\\
US-11 & CRUD особистих дедлайнів & Предмети та дедлайни & Must & Q1 & Sprint 2 & Богдан & 3\\
US-12 & Групові дедлайни (створює староста, бачать усі) & Предмети та дедлайни & Must & Q1 & Sprint 2 & Богдан & 3\\
US-13 & Список дедлайнів з фільтрами й сортуванням & Предмети та дедлайни & Must & Q1 & Sprint 2 & Артемій & 3\\
US-14 & Відмітка виконання дедлайну & Предмети та дедлайни & Must & Q1 & Sprint 2 & Артемій & 1\\
US-15 & Форми створення/редагування дедлайну & Предмети та дедлайни & Must & Q1 & Sprint 2 & Артемій & 2\\
US-16 & Календар на тиждень & Календар & Must & Q1 & Sprint 2 & Артемій & 3\\
US-17 & Календар на місяць & Календар & Should & Q2 & Sprint 3 & Артемій & 2\\
US-18 & Експорт у iCal (.ics) & Календар & Could & Q2 & Після семестру & Богдан & 2\\
US-19 & Канбан-дошка командних задач (To Do / Doing / Done) & Канбан задач & Should & Q2 & Sprint 3 & Артемій & 3\\
US-20 & API задач: CRUD, статуси, виконавець & Канбан задач & Should & Q2 & Sprint 3 & Богдан & 2\\
US-21 & Drag-and-drop карток між колонками & Канбан задач & Could & Q2 & Sprint 3 & Артемій & 2\\
US-22 & Коментарі до задач & Канбан задач & Could & Q4 & Після семестру & Обидва & 2\\
US-23 & E-mail-нагадування за N днів до дедлайну & Нагадування & Should & Q2 & Sprint 3 & Богдан & 3\\
US-24 & Налаштування нагадувань користувачем & Нагадування & Could & Q2 & Sprint 4 & Артемій & 1\\
US-25 & Щотижневий дайджест на e-mail & Нагадування & Could & Q4 & Після семестру & Богдан & 2\\
US-26 & Telegram-бот для нагадувань & Нагадування & Won't (семестр) & Q4 & Після семестру & Богдан & 5\\
US-27 & Прогрес групи по дедлайну (\% виконали) & Аналітика & Should & Q2 & Sprint 3 & Обидва & 2\\
US-28 & Особиста статистика виконання & Аналітика & Could & Q4 & Sprint 4 & Артемій & 2\\
US-29 & Навантаження групи по тижнях (графік) & Аналітика & Could & Q4 & Після семестру & Артемій & 3\\
US-30 & Імпорт дедлайнів з Google Classroom & Інтеграції & Won't (семестр) & Q2 & Після семестру & Богдан & 5\\
US-31 & Синхронізація з Google Calendar & Інтеграції & Won't (семестр) & Q2 & Після семестру & Богдан & 5\\
US-32 & PWA / мобільна версія & Інтеграції & Won't (семестр) & Q4 & Після семестру & Артемій & 5\\
US-33 & Масштабування на факультет: кілька груп, адмін-панель & Інтеграції & Won't (семестр) & Q4 & Після семестру & Обидва & 8\\
D-01 & Звіти за етапами 2–4 у класрум & Інфраструктура та якість & Must & Q3 & Sprint 1 & Обидва & 1\\
D-02 & Щотижневі демо-презентації та рефлексія спринта & Інфраструктура та якість & Must & Q3 & Sprint 2 & Артемій & 1\\
D-03 & Фінальний звіт за модуль 2 і фінальна презентація & Інфраструктура та якість & Must & Q2 & Sprint 4 & Обидва & 3\\
\end{longtable}
\end{landscape}

\section{Обсяг робіт на поточний семестр}
До реалізації в межах модуля~2 відібрано елементи зі спринтів 1--4: усі \textbf{Must} та \textbf{Should}, а також кілька \textbf{Could}, які мають низьку вартість (налаштування нагадувань, особиста статистика, деплой). Разом --- \textbf{30 елементів, 76 story points зі 115} у повному беклозі. Решта (інтеграції, PWA, Telegram-бот, масштабування) залишається в беклозі для версій 1.0--2.5.

\subsection{Спринти та їхні цілі}
\begin{table}[h]
\centering
\renewcommand{\arraystretch}{1.25}
\begin{tabularx}{\textwidth}{l L{2.6cm} Y c}
\toprule
\textbf{Спринт} & \textbf{Період} & \textbf{Ціль спринта} & \textbf{SP}\\
\midrule
Sprint 1 & 28.10 -- 03.11 & Каркас проєкту, CI, Docker; реєстрація та вхід; групи, інвайт-код, ролі & 22\\
Sprint 2 & 04.11 -- 10.11 & Предмети, особисті та групові дедлайни, список із фільтрами, календар тижня; UML і схема БД & 24\\
Sprint 3 & 11.11 -- 18.11 & Канбан задач, e-mail-нагадування, календар місяця, прогрес групи; Swagger, 20+ модульних тестів, E2E & 22\\
Sprint 4 & 19.11 -- 01.12 & Деплой demo-версії, налаштування нагадувань, статистика; фінальний звіт і презентація & 8\\
\midrule
\multicolumn{3}{r}{\textbf{Разом}} & \textbf{76}\\
\bottomrule
\end{tabularx}
\caption{План спринтів}
\end{table}

Навантаження за учасниками: Богдан --- 15 елементів (34 SP), Артемій --- 15 елементів (32 SP), спільні --- 5 елементів (10 SP).

\subsection{Перелік елементів на семестр}
\footnotesize
\renewcommand{\arraystretch}{1.1}
\begin{longtable}{l l L{8.4cm} l c}
\toprule
\textbf{Спринт} & \textbf{ID} & \textbf{Назва} & \textbf{Виконавець} & \textbf{SP}\\
\midrule
\endfirsthead
\toprule
\textbf{Спринт} & \textbf{ID} & \textbf{Назва} & \textbf{Виконавець} & \textbf{SP}\\
\midrule
\endhead
\bottomrule
\endfoot
Sprint 1 & T-01 & Монорепозиторій Next.js + NestJS, лінтери, Prettier & Богдан & 2\\
Sprint 1 & T-02 & Docker Compose: PostgreSQL + backend + frontend & Богдан & 2\\
Sprint 1 & T-03 & CI (GitHub Actions): лінт, збірка, тести & Богдан & 2\\
Sprint 1 & T-04 & Базовий UI-кіт та макет сторінок & Артемій & 3\\
Sprint 3 & T-05 & Swagger / OpenAPI-документація API & Богдан & 2\\
Sprint 3 & T-06 & Модульні тести backend ($\geq$20) & Богдан & 3\\
Sprint 3 & T-07 & Інтеграційні та E2E-тести & Артемій & 3\\
Sprint 2 & T-08 & UML-діаграми та схема БД & Обидва & 2\\
Sprint 4 & T-09 & Деплой на хостинг (Vercel + Railway/Render) & Богдан & 2\\
Sprint 1 & US-01 & Реєстрація за e-mail і паролем & Богдан & 2\\
Sprint 1 & US-02 & Вхід / вихід, JWT-сесія & Богдан & 2\\
Sprint 1 & US-03 & Сторінки реєстрації та входу & Артемій & 2\\
Sprint 2 & US-04 & Профіль користувача (ім'я, група, зміна пароля) & Артемій & 2\\
Sprint 1 & US-06 & Створення навчальної групи & Богдан & 2\\
Sprint 1 & US-07 & Приєднання за інвайт-кодом & Обидва & 2\\
Sprint 1 & US-08 & Ролі: староста / учасник & Богдан & 2\\
Sprint 2 & US-09 & Сторінка групи: учасники, керування & Артемій & 2\\
Sprint 2 & US-10 & CRUD предметів групи & Богдан & 2\\
Sprint 2 & US-11 & CRUD особистих дедлайнів & Богдан & 3\\
Sprint 2 & US-12 & Групові дедлайни (створює староста, бачать усі) & Богдан & 3\\
Sprint 2 & US-13 & Список дедлайнів з фільтрами й сортуванням & Артемій & 3\\
Sprint 2 & US-14 & Відмітка виконання дедлайну & Артемій & 1\\
Sprint 2 & US-15 & Форми створення/редагування дедлайну & Артемій & 2\\
Sprint 2 & US-16 & Календар на тиждень & Артемій & 3\\
Sprint 3 & US-17 & Календар на місяць & Артемій & 2\\
Sprint 3 & US-19 & Канбан-дошка командних задач (To Do / Doing / Done) & Артемій & 3\\
Sprint 3 & US-20 & API задач: CRUD, статуси, виконавець & Богдан & 2\\
Sprint 3 & US-21 & Drag-and-drop карток між колонками & Артемій & 2\\
Sprint 3 & US-23 & E-mail-нагадування за N днів до дедлайну & Богдан & 3\\
Sprint 4 & US-24 & Налаштування нагадувань користувачем & Артемій & 1\\
Sprint 3 & US-27 & Прогрес групи по дедлайну (\% виконали) & Обидва & 2\\
Sprint 4 & US-28 & Особиста статистика виконання & Артемій & 2\\
Sprint 1 & D-01 & Звіти за етапами 2–4 у класрум & Обидва & 1\\
Sprint 2 & D-02 & Щотижневі демо-презентації та рефлексія спринта & Артемій & 1\\
Sprint 4 & D-03 & Фінальний звіт за модуль 2 і фінальна презентація & Обидва & 3\\
\end{longtable}
\normalsize

\subsection{Критерії завершення семестру (Definition of Done модуля 2)}
\begin{itemize}[leftmargin=1.5em,nosep]
  \item усі Must-елементи спринтів 1--3 у статусі Done;
  \item $\geq$20 модульних тестів проходять у CI, є журнал результатів;
  \item інтеграційні/E2E-тести покривають сценарії: реєстрація $\to$ створення групи $\to$ груповий дедлайн $\to$ відмітка виконання;
  \item Swagger-документація доступна за \texttt{/api/docs};
  \item UML (use case, class, sequence) та ER-схема включені у звіт;
  \item demo-версія розгорнута та доступна за публічним посиланням.
\end{itemize}

\section{Результат етапу}
Складено дорожню карту проєкту (Project Overview, діаграма Ґанта на 12 місяців, RACI-матриця на 18 завдань), сформовано беклог із 45 елементів, пріоритезований за MoSCoW та матрицею Ейзенхауера, і визначено обсяг семестру --- 30 елементів (76 SP) у чотирьох спринтах. Беклог і план спринтів синхронізовані з простором проєкту в Notion. Наступний етап --- реалізація прототипу зі щотижневими демо.

\end{document}
